% Part: sets-functions-relations
% Chapter: infinite
% Section: dedekind-induction

\documentclass[../../../include/open-logic-section]{subfiles}

\begin{document}

\olfileid{sfr}{infinite}{induction}
\olsection[अंकगणितीय विगमन]{डेडेकिंड बीजसंरचना आणि अंकगणितीय विगमन}

आता निर्णायक बाब अशी की एखादी डेडेकिंड बीजसंरचना---खरे तर,
\emph{कोणतीही} डेडेकिंड बीजसंरचना---नैसर्गिक संख्यांचे पर्यायी प्रतिरूप
म्हणून उपयोगी पडेल. हे पुढील सहज निष्पत्तीमुळे घडते:

\begin{thm}[अंकगणितीय विगमन]\ollabel{thm:dedinfiniteinduction}
$N, s, o$ यांनी मिळून डेडेकिंड बीजसंरचना बनते असे समजा. मग कोणत्याही
संच $X$ साठी:
	\begin{center}
		जर $o \in X$ आणि $(\forall {n} \in N \cap X){s}({n}) \in X$, {तर} $N \subseteq X$.
	\end{center}
\end{thm}

\begin{proof}
डेडेकिंड बीजसंरचनेच्या व्याख्येनुसार $N = \closureofunder{s}{o}$.
आता ${o} \in X$ आणि $(\forall {n} \in N)(n \in X \lif
{s}({n}) \in X)$ या दोन्ही अटी पूर्ण होत असतील, तर
$N = \closureofunder{s}{o} \subseteq X$.
\end{proof}

विगमन हा नैसर्गिक संख्यांचा वैशिष्ट्यपूर्ण गुणधर्म असल्यामुळे मुद्दा असा
आहे: कोणताही डेडेकिंड-अनंत संच दिला असता, आपण डेडेकिंड बीजसंरचना बनवू
शकतो आणि ती बीजसंरचना नैसर्गिक संख्यांचे पर्यायी प्रतिरूप म्हणून वापरू
शकतो.

हे मान्यच की \olref{thm:dedinfiniteinduction} मध्ये विगमन
\emph{संचसैद्धान्तिक} भाषेत मांडले आहे. पण हे तत्त्व अधिक परिचित वाटेल
अशा भाषेत आपण सहज मांडू शकतो:

\begin{cor}\ollabel{natinductionschema}
$N, s, o$ यांनी मिळून डेडेकिंड बीजसंरचना बनते असे समजा. मग प्राचले असू
शकतील अशा कोणत्याही सूत्र $\phi(x)$ साठी:
\begin{center}
	जर $\phi(o)$ आणि $(\forall {n} \in N)(\phi(n)\lif
	\phi({s}({n})))$, {तर} $(\forall n \in N)\phi(n)$
\end{center}
\end{cor}

\begin{proof}
$X = \Setabs{n \in N}{\phi(n)}$ असे मानून आता
\olref{thm:dedinfiniteinduction} वापरा.
\end{proof}

या निष्कर्षात आपण सूत्राला ``प्राचले असण्याबद्दल'' बोललो. याचा ढोबळ
अर्थ असा की कोणत्याही वस्तू $c_1, \ldots, c_k$ साठी आपण
$\phi(x, c_1, \ldots, c_k)$ बरोबर काम करू शकतो. अधिक नेमकेपणाने,
``प्राचले'' असा उल्लेख न करता निष्कर्ष पुढीलप्रमाणे मांडता येतो.
ज्याची सर्व मुक्त चरे दाखवली आहेत अशा कोणत्याही सूत्र
$\phi(x, v_1, \ldots, v_k)$ साठी:
	\begin{align*}
			\forall v_1 \ldots \forall v_k((&\phi(o, v_1,\ldots, v_k) \land {}\\
			&	(\forall x \in N)(\phi(x,v_1, \ldots, v_k) \lif \phi(s(x), v_1,\ldots, v_k))) \lif {}\\
			&\hspace{3em} (\forall x \in N)\phi(x, v_1,\ldots, v_k))
	\end{align*}
``प्राचले असणे'' असे म्हणाल्याने मजकूर वाचणे फारच सोपे होते, हे उघड आहे.
(\olref[sth][][]{part} मध्ये आपण ही युक्ती वारंवार वापरू.)

पुन्हा डेडेकिंड बीजसंरचनांकडे वळू: कोणतीही डेडेकिंड बीजसंरचना दिली
असता, बेरीज, गुणाकार आणि घातांक क्रिया यांची नेहमीची अंकगणितीय फलनेही
आपण परिभाषित करू शकतो. मात्र हे क्षुल्लक नाही आणि त्यात
\emph{पुनरावर्ती व्याख्येचे} तंत्र वापरावे लागते. हे तंत्र आपण बरेच नंतर,
त्यापेक्षा अधिक व्यापक संदर्भात मांडू आणि समर्थित करू.
(उत्सुक वाचकांनी \olref[sth][ord-arithmetic][]{chap} नंतर याकडे पुन्हा
वळावे किंवा \citealt[pp.~95--8]{Potter2004} यांसारखा पर्यायी विवेचनग्रंथ
वाचावा.) पण $N, s, o$ यांनी डेडेकिंड बीजसंरचना बनत असेल, तर शेवटी
आपल्याला पुढीलप्रमाणे अर्थ ठरवता येईल:
\begin{align*}
	{a} + {o} &= {a} & & & {a} \times {o} &= {o} & & & {a}^{o} &= s(o)\\
{a} + {s}({b}) &= {s}({a}+{b}) &&& {a} \times {s}({b}) &= ({a}\times {b}) + {a}  & & & {a}^{{s}({b})} &= {a}^{b} \times {a}
\end{align*}
आणि ही फलने आपल्या अपेक्षेप्रमाणे वागतात, हे दाखवता येईल.

\end{document}

